\begin{table*}[!htbp]
\centering
\caption[Calibrare SIR pe România]{Calibrare SIR pe România, valul 1 ($15$~martie -- $15$~august 2020 antrenare, $16$~august -- $15$~septembrie 2020 test). Două scenarii: \emph{liber} (optimizează ambele rate, $(\beta,\gamma)$), \emph{constrâns} (folosește $\gamma=1/14$ din durata clinică medie a infecției COVID-19, optimizează doar $\beta$). Comparația ilustrează tensiunea dintre minimizarea pierderii și interpretabilitatea epidemiologică a parametrilor.}
\label{tab:calibration_results}
\begin{tabularx}{\textwidth}{l r r X}
\toprule
\textbf{Marime} & \textbf{Liber} & \textbf{Constrâns} & \textbf{Observatii} \\
\midrule
$\hat{\beta}$ (zi$^{-1}$) & 0.3677 & 0.1039 & rata de transmitere \\
$\hat{\gamma}$ (zi$^{-1}$) & 0.3346 & 0.0714 (\textit{fix}) & $1/\hat{\gamma}$ liber $= 3.0$ zile (non-fizic), constrâns $= 14$ zile \\
$\hat{\mathcal{R}}_0$ & 1.10 & 1.45 & valori de literatură pentru tulpina originala $\sim 2$--$3$ \\
\midrule
MAE antrenare & 1630 & 1612 & cazuri \\
RMSE antrenare & 2199 & 2203 & cazuri \\
MAE test & 13031 & 16727 & extrapolare 1 luna \\
RMSE test & 14889 & 19371 & cazuri \\
\bottomrule
\end{tabularx}
\end{table*}
